\section{Anchors: does the pipeline reproduce what is already known?}
\label{sec:anchors}

Everything that follows depends on the measurement chain being sound, so the
atlas opens with two settings whose answers are known exactly. Neither is a
discovery. Their purpose is that a regression anywhere in the chain --- an
order-parameter normalisation, a censoring rule, a finite-size fit --- shows up
here as a number rather than as a vague loss of agreement downstream.

\subsection{Gardner storage capacity}

A spherical perceptron in $d$ dimensions is asked to separate $P=\alpha d$
random $\pm1$ labels on random Gaussian patterns. The classical result is that
the probability of success tends to one below $\alpha_c=2$ and to zero above
it, with the step sharpening as $d$ grows.

Separability is decided by a linear program rather than by training: a weight
vector satisfying $y_i\,w^\top x_i\ge1$ for all patterns exists if and only if
the labelling is separable. This matters for an anchor. Training to zero error
and declaring failure when it is not reached would conflate ``not separable''
with ``not optimised long enough'', which is exactly the confound the anchor
block exists to exclude.

\manuscriptfigure{figures/family_perceptron_capacity.pdf}{%
  \textbf{The anchor.} Separability probability (A) steps from one to zero near
  the Gardner capacity and the step steepens with $d$; the disorder
  susceptibility (B) peaks at the same load
  $\alpha=\AtlasPerceptronCapacityPeakControl{}$ at every dimension, with
  height growing as
  $N^{\AtlasPerceptronCapacitySharpening{}}$. This is the only family in the
  release to pass every gate --- and its answer was known in advance.%
}{fig:perceptron}

\paragraph{Observation.}
Over \AtlasPerceptronCapacityConditions{} conditions spanning
$\alpha\in\AtlasPerceptronCapacityControlRange{}$ at
\AtlasPerceptronCapacitySizes{} dimensions with
\AtlasPerceptronCapacitySeeds{} outer seeds, the measured separability
probability holds at one well below the predicted capacity, falls through the
neighbourhood of $\alpha=2$, and reaches zero above it. The disorder
susceptibility has an interior maximum at
$\alpha=\AtlasPerceptronCapacityPeakControl{}$ --- against the exact
$\alpha_c=2$, and at the grid point nearest to it --- and that location is
\emph{identical at all four dimensions}, while the peak height grows as
$N^{\AtlasPerceptronCapacitySharpening{}}$ with
$R^2=\AtlasPerceptronCapacitySharpeningRSq{}$. The Binder curves cross at
$\alpha=\AtlasPerceptronCapacityCrossingControl{}$. The verdict is
\emph{\AtlasPerceptronCapacityVerdict{}}, grade
\AtlasPerceptronCapacityGrade{}.

\paragraph{Interpretation.}
The pipeline recovers the transition at the right location, and its finite-size
machinery behaves as it should on a case where the truth is known. The
size-independence of the peak is worth stating explicitly: a critical capacity
that does not move with $d$ is exactly what the Gardner result asserts, and it
is what makes the contrast with \S\ref{sec:symmetry} --- where the peak
\emph{does} move --- a measurement rather than an artefact of the analysis.

This is also the only family in the release that the decision procedure awards
its strongest verdict. We regard that as the main evidence that the procedure
is calibrated rather than permissive: applied to fourteen settings whose
answers are not known, it declined to certify a transition in every one, and
certified the single case that was placed there because the answer was known
in advance.

A solver status other than ``feasible'' or ``provably infeasible'' is recorded
per run as a flag, so no numerical failure can be silently absorbed into the
order parameter.

\paragraph{Not established.}
We do not fit a critical exponent here. The grid was chosen to bracket a known
answer, not to resolve the width of the transition, and the binary order
parameter makes the susceptibility a coarse instrument at these sizes.

\subsection{The Gaussian-mixture error formula}

For an isotropic two-component mixture with centres $\pm\mu v$, a unit rule $w$
at angle $\theta$ to $v$ has generalisation error exactly
\begin{equation}
  \epsilon_g(\theta)=\Phi\!\left(-\mu\cos\theta\right),
  \qquad
  \epsilon_g^{\star}=\Phi(-\mu),
  \label{eq:bayes}
\end{equation}
the second being the Bayes error attained at $\cos\theta=1$. Equation
\eqref{eq:bayes} has no free constants, so comparing it against a measured
error is a direct test of the error bookkeeping rather than of any learning
algorithm.

\paragraph{Observation.}
Across \AtlasMixtureBayesConditions{} conditions and
\AtlasMixtureBayesSizes{} dimensions, the measured misclassification rate of a
fitted linear rule agrees with \eqref{eq:bayes} evaluated at the
\emph{measured} $\cos\theta$ to within probe-set sampling noise; a
representative condition gives measured $0.0129$ against a predicted $0.0135$,
a residual of $-5\times10^{-4}$. The excess over the Bayes floor is
non-negative at every condition, as it must be.

\paragraph{Interpretation.}
The function-space error identity \eqref{eq:epsg} and the classification error
convention used throughout the atlas are correct as implemented. Every
$\epsilon_g$ reported in later sections is produced by the same code path that
reproduces \eqref{eq:bayes} here.

The verdict returned for this family is
\emph{\AtlasMixtureBayesVerdict{}}, and that is the correct answer: there is no
transition in $\mu$ to find. The order parameter moves over
$\AtlasMixtureBayesSemanticRatio{}$ times its interval, so the response is
unambiguous, and the analysis declines to invent a boundary in a setting that
has none. An anchor that is not a transition is a useful second test of the
decision procedure --- a permissive procedure would have found something here.

\paragraph{Not established.}
This anchor validates the estimator, not the learning theory: the fitted rule
is ordinary least squares, and its distance from the Bayes direction is
measured rather than predicted. A stronger anchor would compare $\cos\theta$
itself against its asymptotic prediction, which we do not attempt.

\paragraph{Why only two anchors.}
An anchor is only useful if its closed form is unambiguous and short enough to
state without hidden conventions. We considered adding the replica prediction
for ridge regression and the Saad--Solla plateau, and did not include them in
this release: both require normalisation conventions that would have to be
checked as carefully as the thing they are meant to certify, and a wrong anchor
is worse than none. They are the obvious next additions to this block.
